% Part: first-order-logic
% Chapter: models-theories
% Section: theories

\documentclass[../../../include/open-logic-section]{subfiles}

\begin{document}

% SEGMENT OLP-0170-S01

\olfileid{fol}{mat}{the}

\olsection{முதல்தரக் கோட்பாடுகளின் எடுத்துக்காட்டுகள்}

\begin{ex}
மொழி~$\Lang L_<$ இல் கடு நேரியல் வரிசைகளின் கோட்பாடு பின்வரும்
தொகுப்பால் அடிகோளாக்கப்படுகிறது:
\begin{align*}
\{\quad & \lforall[x][\lnot x < x], \\
& \lforall[x][\lforall[y][((x < y \lor y <
    x) \lor x = y)]], \\
& \lforall[x][\lforall[y][\lforall[z][((x < y
      \land y < z) \lif x < z)]]] \quad \}
\end{align*}
இது நோக்கமாகக் கொண்ட !!{structure}s ஐ முழுமையாகப் பிடிக்கிறது:
ஒவ்வொரு கடு நேரியல் வரிசையும் இந்த அடிகோள் அமைப்பின் மாதிரி; மறுதலையாக,
$R$ ஆனது தொகுப்பு $X$ மீதான நேரியல் வரிசை எனில்,
$\Domain M = X$ மற்றும் $\Assign{<}{M} = R$ கொண்ட கட்டமைப்பு
$\Struct M$ இந்தக் கோட்பாட்டின் மாதிரி.
\end{ex}

\begin{ex}
$\Obj 1$ (!!{constant}), $\cdot$ (இரு-இட !!{function}) ஆகியவற்றைக்
கொண்ட மொழியில் குழுக்களின் கோட்பாடு பின்வருவனவற்றால் அடிகோளாக்கப்படுகிறது:
\begin{align*}
& \lforall[x][\eq[(x \cdot \Obj 1)][x]]\\
& \lforall[x][\lforall[y][\lforall[z][\eq[(x \cdot (y \cdot z))][((x
          \cdot y) \cdot z)]]]]\\
& \lforall[x][\lexists[y][\eq[(x \cdot y)][\Obj 1]]]
\end{align*}
\end{ex}

\begin{ex}
பெயானோ எண்கணிதத்தின் கோட்பாடு, எண்கணித மொழி~$\Lang L_A$ இன் பின்வரும்
வாக்கியங்களால் அடிகோளாக்கப்படுகிறது:
\begin{align*}
& \lforall[x][\lforall[y][(\eq[x'][y'] \lif \eq[x][y])]]\\
& \lforall[x][\eq/[\Obj 0][x']]\\
& \lforall[x][\eq[(x + \Obj 0)][x]]\\
& \lforall[x][\lforall[y][\eq[(x + y')][(x + y)']]]\\
& \lforall[x][\eq[(x \times \Obj 0)][\Obj 0]]\\
& \lforall[x][\lforall[y][\eq[(x \times y')][((x \times y) + x)]]]\\
& \lforall[x][\lforall[y][(x < y \liff \lexists[z][\eq[(z' + x)][y])]]]\\
\intertext{மேலும் பின்வரும் வடிவிலான எல்லா வாக்கியங்களும்}
& (!A(\Obj 0) \land \lforall[x][(!A(x) \lif !A(x'))]) \lif \lforall[x][!A(x)]
\end{align*}
கடைசி வடிவத்தில் முடிவிலா எண்ணிக்கையிலான வாக்கியங்கள் இருப்பதால்,
இந்த அடிகோள் அமைப்பு முடிவிலியானது. கடைசி வடிவம்
\emph{தொகுத்தறிதல் திட்டம்} எனப்படுகிறது. (உண்மையில், தொகுத்தறிதல்
திட்டம் இங்கே கூறியதைவிடச் சற்று சிக்கலானது.)

கடைசி அடிகோள்~$<$ இன் \emph{வெளிப்படையான வரையறை}.
\end{ex}

\begin{ex}
தூய கணங்களின் கோட்பாடு கணிதத்தின் அடித்தளங்களில் (மற்றும் அதன்
மெய்யியலில்) முக்கியப் பங்கு வகிக்கிறது. ஒரு கணத்தின் எல்லா
!!{element}s உம் தூய கணங்களாக இருந்தால் அது தூயது. எனவே வெற்றுக் கணம்
தூயதாகும்; ஆனால் கணமாக இல்லாத ஒன்று !!a{element} ஆக உள்ள கணம் தூயதல்ல.
ஆகவே தூய கணங்கள், வெற்றுக் கணத்திலிருந்து மட்டும், தாமே கணங்களாக
இல்லாத பொருட்களான ``அடிப்படையுறுப்புகள்'' எதுவுமின்றி உருவாக்கப்பட்டவை.

பின்வருவதைத் தூய கணங்களின் கோட்பாட்டுக்கான அடிகோள் அமைப்பாகக்
கருதலாம்:
\begin{align*}
& \lexists[x][\lnot \lexists[y][y \in x]]\\
& \lforall[x][\lforall[y][(\lforall[z](z \in x \liff z \in y) \lif 
\eq[x][y])]]\\
& \lforall[x][\lforall[y][\lexists[z][\lforall[u][(u \in z \liff
(\eq[u][x] \lor \eq[u][y]))]]]]\\
& \lforall[x][\lexists[y][\lforall[z][(z \in y \liff \lexists[u][(z \in
        u \land u \in x)])]]]\\
\intertext{மேலும் பின்வரும் வடிவிலான எல்லா வாக்கியங்களும்} &
\lexists[x][\lforall[y][(y \in x \liff !A(y))]]
\end{align*}
முதல் அடிகோள், !!{element}s எதுவுமில்லாத ஒரு கணம் உள்ளது (அதாவது,
$\emptyset$ உள்ளது) என்கிறது. இரண்டாவது, கணங்கள் விரிவுநிலைப்
பண்புடையவை என்கிறது. மூன்றாவது, எந்தக் கணங்கள் $X$ மற்றும் $Y$ க்கும்
கணம் $\{X, Y\}$ உள்ளது என்கிறது. நான்காவது, எந்தக் கணம் $X$ க்கும்
கணம் $\cup X$ உள்ளது என்கிறது; இங்கு $\cup X$ என்பது $X$ இன் எல்லா
உறுப்புகளின் ஒன்றியம்.

கடைசியாகக் குறிப்பிடப்பட்ட !!{sentence}s அனைத்தும் சேர்ந்து
\emph{அப்பாவி உட்கொள்ளல் திட்டம்} எனப்படுகின்றன. சாரமாக, ஒவ்வொரு
$!A(x)$ க்கும் கணம் $\Setabs{x}{!A(x)}$ உள்ளது என்று அது கூறுகிறது;
எனவே முதல் பார்வையில் அது மெய்யான, பயனுள்ள, ஒருவேளை அவசியமான
அடிகோளாகவும் தோன்றுகிறது. ஆனால் அது இக்கோட்பாட்டை
நிறைவுசெய்யத்தக்கதல்லாததாக ஆக்குவதால் ``அப்பாவி'' எனப்படுகிறது:
$!A(y)$ ஐ $\lnot y \in y$ ஆக எடுத்தால், பின்வரும் !!{sentence} கிடைக்கும்:
\[
\lexists[x][\lforall[y][(y \in x \liff \lnot y \in y)]]
\]
இந்த !!{sentence} எந்த !!{structure} இலும் நிறைவுறுத்தப்படாது.
\end{ex}

\begin{ex}
\emph{பகுதியியல்} துறையில், \emph{பகுதியாக இருத்தல்} என்ற தொடர்பு ஓர்
அடிப்படைத் தொடர்பாகும். கணங்களின் கோட்பாடுகளைப் போலவே, இத்தொடர்பைப்
பற்றிய பல்வேறு (சிலவேளைகளில் முரண்படும்) கருத்தாக்கங்களை
அடிகோளாக்கும் பகுதியியல் கோட்பாடுகள் உள்ளன.

பகுதியியலின் மொழி ஒரே இரு-இடப் பயனிலைக் குறியீடு~$\Obj P$ ஐக்
கொண்டுள்ளது; $\Atom{\Obj P}{x, y}$ என்பது $x$,~$y$ இன் ``ஒரு பகுதி''
என்று ``குறிக்கிறது''. இந்தப் பொருள்கோளை மனத்தில் கொண்டிருக்கும்போது,
இம்மொழிக்கான !!a{structure} ஒரு \emph{பகுதியியல் கட்டமைப்பு} எனப்படுகிறது.
ஒரே இரு-இடப் பயனிலைக்கான ஒவ்வொரு கட்டமைப்பும் உண்மையில் இந்தப் பெயருக்குத்
தகுதியாகாது. ``பகுதியாக இருத்தலை'' பிடிக்க வேண்டுமெனில்,
$\Assign{\Obj P}{M}$ சில நிபந்தனைகளை நிறைவுறுத்த வேண்டும்; அவற்றைப்
பகுதியியல் கோட்பாட்டின் அடிகோள்களாக விதிக்கலாம். எடுத்துக்காட்டாக,
பொருட்கள் மீதான பகுதியாக இருத்தல் ஒரு பகுதி வரிசை: ஒவ்வொரு பொருளும்
தன்னுடைய ஒரு பகுதி (என்றாலும் \emph{முறையற்ற} பகுதி); வேறுபட்ட இரண்டு
பொருட்கள் ஒன்றுக்கொன்று பகுதிகளாக இருக்க முடியாது; ஒரு பொருளின் ஒரு
பகுதியின் பகுதியும் அந்தப் பொருளின் பகுதியே. இவ்வுணர்வில் ``ஒரு பகுதி''
என்பது ``ஓர் உட்கணம்'' என்பதை ஒத்திருக்கிறது; ஆனால் தற்சுட்டோ
கடப்போ அல்லாத ``ஓர் உறுப்பு'' என்பதை ஒத்திருக்கவில்லை.
\begin{align*}
& \lforall[x][\Atom{\Obj P}{x,x}] \\
& \lforall[x][\lforall[y][((\Part{x}{y} \land \Part{y}{x})
      \lif \eq[x][y])]] \\
& \lforall[x][\lforall[y][\lforall[z][((\Part{x}{y} \land
        \Part{y}{z}) \lif \Part{x}{z})]]]\\
\intertext{மேலும், எந்த இரண்டு பொருட்களுக்கும் ஒரு பகுதியியல் கூட்டுத்தொகை
  உள்ளது (இரு பொருட்களையும் பகுதிகளாகக் கொண்டு, இவ்வகையில் மிகச்சிறியதாக
  இருக்கும் ஒரு பொருள்).}  &
\lforall[x][\lforall[y][\lexists[z][\lforall[u][(\Part{z}{u} \liff
        (\Part{x}{u} \land \Part{y}{u}))]]]]
\end{align*}
இவை மெய்யியலாளர்கள் கருதும் பகுதியாக இருத்தலின் அடிப்படைக் கொள்கைகளுள்
சில மட்டுமே. ஆனால் வரையறுக்கப்பட்ட சில தொடர்புகளை முதலில்
அறிமுகப்படுத்தாமல் மேலதிகக் கொள்கைகளை உருவாக்கவோ எழுதவோ விரைவிலேயே
கடினமாகிறது. எடுத்துக்காட்டாக, பகுதியியலில் ஆர்வமுள்ள பெரும்பாலான
மெய்யியலாளர்கள் பின்வருவதையும் செல்லுபடியான கொள்கையாகக் கருதுகின்றனர்:
ஒரு பொருள்~$x$ க்கு முறையான பகுதி~$y$ இருக்கும்போதெல்லாம்,~$y$ உடன்
பொதுவான பகுதிகள் எதுவுமில்லாத பகுதி~$z$ உம் அதற்கு உண்டு; மேலும் $y$
மற்றும் $z$ இன் இணைவு~$x$ ஆகும்.
\end{ex}

\end{document}
